\clearpage
\section{Rectangular observations with fewer than $D$ variable coordinates}
\label{sec:partial188}
The one-coordinate dependence in the laminate theorem can be weakened to
arbitrary dependence on any $k<D$ coordinates. The construction below uses
the complete finite Sobolev proof already included in this manuscript.
It retains at least one exactly invariant coordinate. No local spectral
estimate or successful partition is added as an unproved input.

Put $p=D/2$, fix $1\le k<D$, and set $d=D-k$. In this section define
\[
 \rho=(2/3)^d,\qquad C_D=4A_DD^{1/p},\qquad
 F_{D,k}=\frac{2^D3^d}{1-\rho},
\]
where $A_D$ is the explicit constant in
\eqref{eq:rectangle-unweighted-derived}. These constants are not optimal.
\begin{theorem}\label{thm:partial188}
Let $D\ge3$, $n\ge2$, $c>0$, and let
$a(x)=\alpha(x_1,\ldots,x_k)>0$ on $(\mathbb Z/n\mathbb Z)^D$.
There is a disjoint axis-aligned integer rectangular partition $\mathcal P$,
independent of the field, such that for every scalar real or complex field
\[
 c\|u\|^2\le\E_a(u)+c\|C_{\mathcal P}u\|^2,\qquad
 C_{\mathcal P}u=(\sqrt{|R|}\,\overline u_R)_{R\in\mathcal P},
\]
and
\begin{equation}\label{eq:partial-count188}
 \#\mathcal P\le2^d+F_{D,k}(cC_D)^p\sum_Xa^{-p}.
\end{equation}
If $\mu=\operatorname{mean}(a)\ge1$ and
$\operatorname{mean}((\mu/a)^p)\le M$, then for an $h$-component field the
output dimension is at most
\[
 h\bigl[2^d+F_{D,k}(cC_D)^p2^pM\bigr]
 \max\{1,n^D/(1+\mu)^p\}.
\]
The energy inequality holds coordinatewise on all fields, and hence on an
actual coexact subspace. Output dimensions are over the selected scalar
field; double them for complex fields counted over the reals.
\end{theorem}

\subsection{Local energy derived from the original coefficients}
For a cut rectangle $R$ with volume $V$, longest side $L$ and
$\kappa=L^2/V^{2/D}$, the established finite inequality
\eqref{eq:rectangle-weighted-derived} gives
\begin{equation}\label{eq:partial-local188}
 \operatorname{Var}_R(u)\le A_D D^{1/p}\kappa
      \Bigl(\sum_Ra^{-p}\Bigr)^{1/p}\E_{a,R}(u).
\end{equation}
Indeed H\"older uses $q=2D/(D+2)$ and $q/(2-q)=p$, and each original
starting vertex contributes at most $D$ internal edges. Reflection has
already been completed in the unweighted proof; no reflected weights,
enlarged-region energies or assumed Poincar\'e constant enter here.
A singleton has zero variance. Internal cut edges are a subset of the
actual periodic positive starts, including when $n=2$; closing and
cross-cell contributions may be discarded as nonnegative.

\subsection{A balanced base tree and actual extrusion}
The base domain is $\{0,\ldots,n-1\}^k$. For a base rectangle $Q$, let
$m(Q)$ be its longest integer side and $b(Q)=\sum_Q\alpha^{-p}$.
Set $\theta=(cC_D)^{-p}$. Split a node precisely when
$m\ge2$ and $m^d b(Q)>\theta$. Each non-singleton coordinate is bisected
into consecutive floor and ceiling halves; singleton coordinates stay
unchanged. Accept equality and never split a base singleton.

All side lengths at a node differ by at most one, by induction from the
root cube. Bisecting $m$ and $m-1$ preserves this property, also for mixed
lengths one and two. At least one parent side has length $m$, so every
child has longest side $m'$ satisfying
\begin{equation}\label{eq:partial-child188}
 m/3\le\lfloor m/2\rfloor\le m'\le\lceil m/2\rceil\le2m/3
 \qquad(m\ge2).
\end{equation}
The tree terminates, has at most $2^k$ children per internal node, and
all base sides at a node with $m\ge2$ are at least $m/2$.

For each leaf $Q$, partition each of the $d$ invariant coordinates into
$r=\lceil n/m\rceil$ consecutive balanced intervals of lengths
$\lfloor n/r\rfloor$ or $\lceil n/r\rceil$. Every such length lies in
$[m/2,m]$. The upper bound follows from $n\le rm$. For $r\ge2$,
$n\ge(r-1)m+1$ implies
$\lfloor n/r\rfloor\ge\lfloor(m+1)/2\rfloor\ge m/2$; for $r=1$,
$m=n$. Thus every full rectangle $R=Q\times Z$ has $\kappa\le4$.
These rectangles partition the actual torus representatives, with no
wrapping or overlap. If $m=1$, every resulting full rectangle is a singleton.

Exact invariance, rather than a sampled favourable fibre, gives
\[
 \sum_Ra^{-p}=|Z|b(Q)\le m^d b(Q).
\]
For a non-singleton leaf, \eqref{eq:partial-local188} therefore yields
$c\operatorname{Var}_R(u)\le\E_{a,R}(u)$. Singleton leaves need no mass
acceptance condition, since their residual variance is zero. Summing and
using orthogonal mean decomposition proves the actual inequality with
$K=c$ in Theorem~\ref{thm:partial188} for all fields simultaneously.

\subsection{Terminal charging and the critical rank count}
Write $w(Q)=m(Q)^{-d}$. Every internal node satisfies
\[
 w(Q)<\theta^{-1}b(Q).
\]
Let $T$ be internal nodes with no internal children. These nodes are
pairwise disjoint in the base lattice. Assign each internal node to any
terminal internal descendant. Nodes assigned to a fixed terminal $J$ lie
on its ancestor path at distinct depths. By \eqref{eq:partial-child188},
ancestor weights decay by at least the factor $\rho=(2/3)^d<1$, whence
\begin{equation}\label{eq:partial-charge188}
 \sum_{Q\text{ internal}}w(Q)
 \le\frac{1}{1-\rho}\sum_{J\in T}w(J)
 \le\frac{\theta^{-1}}{1-\rho}\sum_{\text{base}}\alpha^{-p}.
\end{equation}
Only terminal internal nodes are charged as disjoint; all internal nodes
are not incorrectly treated as disjoint.

Every child has weight at most $3^d$ times its parent's weight. Exact
leaf-weight telescoping and the arity bound give
\[
 \sum_{Q\text{ leaf}}w(Q)
 \le n^{-d}+2^k3^d\sum_{Q\text{ internal}}w(Q).
\]
A base leaf creates exactly $\lceil n/m\rceil^d\le(2n/m)^d$ full cells.
Multiplying by $(2n)^d$, applying \eqref{eq:partial-charge188}, and using
$n^d\sum_{\text{base}}\alpha^{-p}=\sum_Xa^{-p}$ proves
\eqref{eq:partial-count188}. The accepted root case has one cell and no
internal charges; singleton leaves remain in the telescoping identity.
Finally the moment condition gives
$\sum_Xa^{-p}\le Mn^D/\mu^p\le2^pMn^D/(1+\mu)^p$, proving the stated
output budget. This completes the analytical proof.

\subsection{Verification, attribution and remaining boundary}
The exact source188 proof is pinned at
\texttt{0178b02bf26381633726edb82758e1c7ccb9aefd}; its complete involved-model
source audit and all limitations are retained under
\path{papers/flat-coarse-rank/research/evidence188/} in the source map above.
The audit found no blocking mathematical gap. It did not perform a formal
rebuild or scientific replay. Separate finite tree, geometry, inverse-mass
and rank diagnostics were then executed at
\texttt{b9d9dfdee043a62ec5d30926e215f8a581507b70}: two fresh own Linux Colab
clones, each in normal and optimized modes, produced four identical payloads.
Each checked 648 rational fixtures and rejected six actual mutated controls.
All twelve child exit files contain zero. These checks are not local spectral
computations, all-size proof or independent research replication; their
coefficients are restricted to squares to keep inverse powers rational.
The all-size claim rests on the analytical argument above.

Weighted Sobolev and adaptive-tree mechanisms are classical. The targeted
primary comparison in source187--188 includes Cruz-Uribe, Dal and Rodney,
\emph{Degenerate Sobolev and Poincar\'e inequalities via extrapolation},
arXiv:2608.20498v1, Theorem 4.5
(\url{https://arxiv.org/html/2608.20498v1}), and Claros, P\'erez and Zheng,
\emph{From generalized Poincar\'e to Poincar\'e--Sobolev inequalities via
self-improving methods}, arXiv:2606.08556v1, Appendix C
(\url{https://arxiv.org/html/2606.08556v1}). Their inspected continuum and
summability statements are context, not imported premises or a priority
certificate. Broader approximation literature still requires comparison.

\paragraph{Model assessments and review context.}
The source audit used the historically involved OpenAI Codex agent with a
configured gpt-6-astra context. Effective served backend and reasoning setting
were unexposed; shared history, workspace, account and provider context do not
establish blindness or independence. Its full report, L1--L5 limitations and
factual inspection scope remain preserved. Neither that source audit nor the
full revision187 PDF assessment supplies a grade for this changed PDF. This
complete artifact requires its own numerical overall and criterion stars,
with all objections retained; no human peer review or editorial acceptance
is inferred.

At $k=D$, the argument has $d=0$ and $\rho=1$, so the ancestor sum in
\eqref{eq:partial-charge188} is no longer controlled. This is a limitation
of the present method, not a counterexample to general rectangular means.
Approximate invariance and arbitrary real-coefficient implementation or
numerical stability are not established. Constants are nonoptimal upper-order
bounds, not minimum-rank claims. The sole unrestricted critical rectangular
question remains open, and physical interacting/continuum questions stay separate.